\section{Vendor-backed support and scope discipline}\label{sec:vendor-support}

The theorem package of the present paper is not proved from vendor-backed
evidence. It is proved from the internal structural development and lifted
conditionally through an explicit external contract. The role of the
vendor-backed PICA evidence is different. It constrains the paper's
primitive-role assignment, supports its scope claims, and blocks simpler but
misleading narratives in which one-dimensional scoring or route sensitivity is
allowed to do explanatory work that the theorem package itself does not
support. This section therefore records what the vendor-backed evidence is
doing, and just as importantly, what it is not doing.

\subsection*{What the vendor-backed evidence is doing}

The vendor-backed evidence fixes the empirical perimeter within which the
theorem route is being read. The vendor-backed PICA artifacts show that richer
primitive interaction structure is real, that baseline-like configurations are
not the whole story, and that distinct cells and program families contribute
unevenly. In that sense the evidence supports the claim that the primitive roles
must be assigned deliberately rather than flattened into a single
undifferentiated closure narrative.

More specifically, the vendor-backed analysis supports three narrow claims used
by the paper.

First, baseline-like configurations are weaker than richer structured
configurations when the actual PICA interaction space is examined. That result
supports the decision not to identify theory growth with a bare fixed-package
closure story.

Second, cell importance is heterogeneous. The vendor-backed ablation surfaces do
not behave as though every primitive interaction contributes equally or
interchangeably. That supports the paper's asymmetric primitive-role assignment:
the theorem track is allowed to privilege \(P5\), \(P1/P2\), \(P6\), and \(P4\) without
pretending that all primitive roles collapse to the same explanatory type.

Third, richer interaction structure matters. The vendor-backed analysis rules
out the idea that the present paper could safely replace primitive interaction
with a one-number proxy and lose nothing of consequence. This supports the
paper's narrower theorem route: the broader interaction space is real and not
automatically compressible into the same theorem-bearing language.

\subsection*{What survived from toy to vendor contact}

The paper's early toy analysis served only as a negative filter on naive
abstractions. In small finite spaces, broad uniqueness and broad stability
claims fail quickly under simplistic frontier semantics. That negative lesson
survives vendor contact in a limited but important sense: the paper still does
not earn a broad empirical or theorem-level canonicality claim merely by
introducing an efficiency language.

At the same time, the vendor-backed evidence improves the structural picture in
ways the toy analysis could not. The toy phase was useful for falsifying naive
claims, but it was too coarse to ground any lawfulness story in the actual PICA
regime. Vendor-backed analysis at the shipped scales supports a more disciplined
conclusion: structured configurations matter, cell importance is uneven, and the
interaction space should not be treated as a trivial decoration on top of
closure.

What survives across both phases is therefore a negative-to-positive division of
labor. The toy phase justifies caution against overbroad frontier rhetoric. The
vendor-backed phase justifies the paper's narrower role assignment and helps
keep proof language from outrunning interaction structure. Together they
support the paper's scope limits: a theorem-led conditional paper rather than a
broad canonical-collapse claim.

\begin{table}[t]
\caption{Boundary between toy falsification and vendor-backed support.}
\label{tab:toy-vendor-boundary}
\centering
\small
\begin{tabularx}{\textwidth}{@{}l X X@{}}
\toprule
Question & Toy phase & Vendor-backed analysis \\
\midrule
Broad uniqueness/stability & naive uniqueness and broad stability fail quickly under simplistic frontier semantics & such failures remain a boundary condition; no vendor evidence upgrades them into broad canonicality \\[4pt]
Baseline-only closure narrative & too coarse to justify lawfulness claims & baseline-like configurations are weaker than richer structured configurations \\[4pt]
Homogeneous primitive importance & not justified & cell importance is heterogeneous across the shipped analysis regime \\[4pt]
One-number interaction summary & too coarse for a realistic interaction story & richer interaction structure matters and should not be collapsed into a single proxy \\[4pt]
Main conditional theorem & not supported & supporting-only scope discipline; no theorem proof support \\
\bottomrule
\end{tabularx}
\end{table}

\subsection*{What the vendor-backed evidence is not doing}

The vendor-backed evidence is not theorem support. It does not discharge the
fixed-package saturation theorem. It does not prove the package-change bridge
theorem. It does not prove the restricted alignment lemma. And it does not
discharge the external arithmetic assumptions in the conditional lift. None of
those proof obligations are delegated to experiment.

This boundary must remain explicit because the paper uses vendor evidence near
the theorem route without allowing it to cross the proof boundary. The internal
theorem package is formalized in Lean. The arithmetic lift depends on named
external assumptions recorded in the assumption box and external dependency
contract. The vendor-backed evidence belongs to neither of those proof layers.
It is supporting-only.

Nor does the vendor-backed evidence justify stronger claims than the paper
actually makes. It does not establish unrestricted global canonicality. It does
not show that every evaluation regime must select the same representative
family. It does not erase the need for the explicit external arithmetic
contract. And it does not license the use of \(P3\) as a theorem-driving
primitive merely because route sensitivity is visible in richer interaction
data. The vendor-backed material is used to support scope discipline, not to
extend the theorem's conclusion.

\subsection*{Scope discipline}

The right lesson of the vendor-backed evidence is not that theorem work should
be abandoned. It is that theorem work should stay within a carefully frozen
slice of the larger interaction space. That is exactly what the present paper
does. \(P5\) remains the fixed-package object, \(P1/P2\) remain the mechanisms of
package change, \(P6\) remains the frozen evaluation regime, \(P4\) remains the
support/index axis, and \(P3\) remains diagnostic-only. The vendor-backed
evidence supports this cleanup because it shows that broader interaction
structure exists and must not be casually collapsed.

This also explains why the present paper is theorem-led but conditional. The
theorem route carries a central structural claim within a sharply delimited
setting. The empirical side is correspondingly
supportive rather than constitutive. It constrains narrative inflation,
clarifies which primitive roles are defensible, and supports the paper's
out-of-scope declarations. Section~\ref{sec:scope-limits} makes those limits
explicit.
